\section{From a weak Hessian bound to variable curvature}
\label{sec:curvature-change}

We now apply the generator facts to \(m_\lambda=e^{\lambda F}m\).
The weak Hessian upper bound lowers the Ricci bound by
\(\lambda G\). We first prove this identity on compact original
test functions, then close it in the common generator domain. This
order is needed because a Lipschitz drift can destroy the Sobolev
regularity of an original test function's weighted Laplacian.

For a probability measure \(\mu\) with bounded support, define
\[
\Ent_m(\mu)=
\begin{cases}
\displaystyle\int_M\rho\log\rho\dd m,&\mu=\rho m,\\
+\infty,&\mu\not\ll m,
\end{cases}
\qquad 0\log0=0.
\]
The negative part is integrable: a bounded ball containing the
support has finite measure, and \(s\log s\ge-1/e\).
This is the entropy domain used below, also for the bounded
reweightings \(m_\lambda\).

For a bounded continuous function \(k\), we use the variable
curvature-dimension condition \(\CD(k,\infty)\) of
\cite[Definition~1.2]{BHS2021}. In the bounded-support case needed
below, it supplies, for any two finite-entropy probability measures
\(\mu_0,\mu_1\), a dynamical optimal plan \(\pi\) with
\(\mu_t=(e_t)_\#\pi\) such that
\begin{align}
\label{weight:variable-entropy}
\Ent_m(\mu_t)
&\le (1-t)\Ent_m(\mu_0)+t\Ent_m(\mu_1)\notag\\
&\quad-\int_{\Geo(M)}\int_0^1
  \bigl(\min\{s,t\}-st\bigr)k(\gamma_s)|\dot\gamma|^2
  \dd s\dd\pi(\gamma).
\end{align}
Here \(\Geo(M)\) is the space of constant-speed minimizing paths
on \([0,1]\), and \(e_t(\gamma)=\gamma_t\). A dynamical optimal
plan is a probability measure on this path space whose endpoint
coupling is optimal for quadratic transport. This formula fixes
our Green-kernel and Laplacian normalizations.

\begin{proposition}[A distributional Hessian bound changes curvature]
\label{weight:variable-cd}
Let \((M,d,m)\) be a full-support \(\RCD(K,N)\) space,
where \(K\in\mathbb R\) and \(1<N<\infty\).
Let \(F\colon M\to\mathbb R\) be bounded and globally Lipschitz, and let
\(G\colon M\to\mathbb R\) be bounded and continuous. Suppose
that for every compactly supported \(g\in\Test\) and every
nonnegative \(h\in\Lip_c(M)\),
\[
H_F(\nabla g,\nabla g)(h)
\le\int hG\Gamma(g)\dd m.
\]
For every \(\lambda>0\), the measure
\(m_\lambda=e^{\lambda F}m\) satisfies
\(\CD(K-\lambda G,\infty)\). The metric is unchanged.
\end{proposition}

\begin{proof}
\textbf{Core identity.}
Fix \(\lambda>0\), and write \(V=\lambda F\), \(w=e^V\),
\(k=K-\lambda G\). Take a nonnegative \(\xi\) in
\eqref{weight:multiplier-class}. Lemma~\ref{weight:resolvent}
makes \(\xi\) Lipschitz. Initially let \(g\in\Test\) be
compactly supported, and set
\(q=\Gamma(g)/2\), \(b=\Gamma(V,g)\). Define
\begin{equation}
\label{weight:integrated-bochner}
\mathcal B_V(g,\xi)
=-\int w\Gamma(\xi,q)\dd m
 +\int w\bigl[\xi\Delta_Vg+\Gamma(\xi,g)\bigr]
                  \Delta_Vg\dd m.
\end{equation}
The notation \(\mathcal B_0(g,\eta)\) denotes the same formula
with \(V=0\); it is also meaningful for a bounded Lipschitz
multiplier \(\eta\) after localization near \(\supp g\).
All terms are integrable: \(q\in W^{1,2}\),
\(\Delta_Vg\in L^2\), and \(\xi,|\nabla\xi|\) are bounded.

The original measure-valued Bochner inequality, integrated against
\(w\xi\), gives
\begin{equation}
\label{weight:original-bochner}
\mathcal B_0(g,w\xi)
\ge K\int w\xi\Gamma(g)\dd m.
\end{equation}
Here and below a compact test cutoff equal to one near
\(\supp g\) can be inserted in the multiplier. Its derivatives
do not meet the support of any term. Thus the bounded Lipschitz
multiplier is admissible by the local Bochner-measure calculus
in Section~\ref{sec:preliminaries}.

The product rule and \eqref{weight:generator} give the exact identity
\begin{align}
\label{weight:bochner-difference}
\mathcal B_V(g,\xi)-\mathcal B_0(g,w\xi)
 &=\int w\xi\Gamma(V,q)\dd m
   +\int wb\bigl[\Gamma(\xi,g)+\xi\Delta_Vg\bigr]\dd m\nonumber\\
 &=-H_V(\nabla g,\nabla g)(w\xi).
\end{align}
The second equality is precisely \eqref{eq:weak-hessian}, since
\(\divg(w\xi\nabla g)=w[\Gamma(\xi,g)+\xi\Delta_Vg]\).
Use \eqref{weight:hessian-assumption}, with the same harmless
cutoff, in \eqref{weight:bochner-difference}. Together with
\eqref{weight:original-bochner} this yields
\begin{equation}
\label{weight:core-be}
\mathcal B_V(g,\xi)\ge\int k\xi\Gamma(g)\dd m_V.
\end{equation}

The example in Remark~\ref{weight:domain-warning} explains why
the weighted generator domain needs a separate closure argument.

\smallskip\noindent\textbf{Closure in the generator domain.}
We next close this inequality in the correct generator domain.
Integration by parts in its first term rewrites the left side as
\begin{equation}
\label{weight:graph-form}
\frac12\int (\Delta_V\xi)\Gamma(g)\dd m_V
 +\int (\Delta_Vg)\Gamma(\xi,g)\dd m_V
 +\int\xi(\Delta_Vg)^2\dd m_V.
\end{equation}
Every term in \eqref{weight:graph-form}, as well as the right
side of \eqref{weight:core-be}, is continuous under convergence
in the \(\Delta_V\) graph norm. Indeed, graph convergence implies
strong \(L^2\) convergence of gradients; use the boundedness of
\(\xi,\Delta_V\xi,|\nabla\xi|\), and \(k\).
By Lemma~\ref{pre:approximation} and graph-norm equivalence,
compact original test functions form a core for \(\Delta_V\).
Consequently \eqref{weight:core-be} in the form
\eqref{weight:graph-form} holds for every \(g\in D(\Delta_V)\).
If in addition \(\Delta_Vg\in W^{1,2}\), the last two terms
in \eqref{weight:graph-form} combine by integration by parts,
and give
\begin{equation}
\label{weight:full-be}
\frac12\int (\Delta_V\xi)\Gamma(g)\dd m_V
-\int\xi\Gamma(g,\Delta_Vg)\dd m_V
\ge\int k\xi\Gamma(g)\dd m_V.
\end{equation}
This is the weighted Bakry--\'Emery inequality on its actual
test domain. In particular it holds on the smaller class also
requiring \(\Gamma(g)\in L^\infty\).

\smallskip\noindent\textbf{Constant curvature.}
We first obtain a weighted \(\RCD\) structure with the same
intrinsic distance \(d\). Replace \(k\) by the constant
\(K_V=\inf_M k> -\infty\), and set
\(\mathcal E_V=2\mathrm{Ch}_V\). The unchanged weak gradients
make \(\mathcal E_V\) a strongly local symmetric Dirichlet form
with carr\'e du champ \(\Gamma\).
In the metric topology \(\tau\), using the completed Borel
\(\sigma\)-algebra, completeness, separability, full support
and finite bounded-ball measures give condition \textup{(MD)}
of \cite{AGS2015}.
The original Sobolev-to-Lipschitz property gives \textup{(ED.a)};
boundedly supported Lipschitz functions with slope at most one
have weak gradient at most one, giving \textup{(ED.b)}.
The converse of \cite[Theorem~3.9]{AGS2015} therefore makes
\((M,\tau,m_V,\mathcal E_V)\) an Energy measure space and
identifies its intrinsic distance with \(d\).
Since \(\mathcal E_V=2\mathrm{Ch}_V\) for this metric,
the necessity direction of \cite[Theorem~3.14]{AGS2015}
gives upper regularity. Bounded equivalence preserves Gaussian
volume growth, hence \textup{(MD.exp)}, so the final assertion
of that theorem gives conservativity.
Sobolev-to-Lipschitz also supplies continuous representatives
for energy functions with \(\Gamma\le1\), as required by
\cite[Definition~3.16]{AGS2015}. Thus this is a Riemannian
Energy measure space. Applying \cite[Theorem~4.17]{AGS2015}
to \eqref{weight:full-be} with the constant \(K_V\) gives
\(\RCD(K_V,\infty)\).

\smallskip\noindent\textbf{Variable curvature.}
On this weighted \(\RCD\) space, \eqref{weight:full-be} is its
\(\mathrm{BE}_2(k,\infty)\) condition, and the bounded continuous
function \(k\) meets its lower semicontinuity and lower-bound
assumptions. Hence \(\CD(k,\infty)\) follows
\cite[Definition~1.4 and Theorem~1.1]{BHS2021}.
\end{proof}

\begin{remark}[Why the domain argument is needed]
\label{weight:domain-warning}
An original test function need not satisfy
\(\Delta_Vg\in W^{1,2}\). Even in Euclidean space, a bounded
Lipschitz \(V\) agreeing with \(-|x_1|\) near the origin and a
smooth compactly supported \(g\) agreeing with \(x_1\) there
give \(\Delta_Vg=-\operatorname{sign}(x_1)\) locally.
The integrated identity \eqref{weight:bochner-difference} and
closure through \eqref{weight:graph-form} avoid differentiating
this drift term. They also explain why no general converse from
an unspecified weak Hessian to geodesic convexity is being used;
compare \cite[Remark~2.4]{BrenaGigli2024}.
\end{remark}

